\section{Introduction}
A sprinkler system imitates rain: water under pressure leaves through a nozzle, breaks into drops and falls on the crop. It can irrigate land that is not level, it can apply small depths, and it can be automated. Its drawbacks are energy (the pressure), wind drift and evaporation from the drops, and wetting of the leaves. This chapter gives the design from the nozzle to the lateral.

\section{Types of sprinkler systems}
Sprinkler systems are \emph{set} (portable or solid-set laterals moved between positions), \emph{continuous-move} (centre pivot, linear move, travelling gun) and \emph{permanent}. Solid-set systems need many sprinklers and pipes but little labour; pivots need little pipe but a smooth ground and an even application rate that rises towards the end.

\section{Nozzle discharge and wetted radius}
The discharge of a nozzle follows the orifice law $q=C_dA\sqrt{2gH}$. For $d=4.8$ mm and a pressure of 350 kPa with $C_d=0.95$, $q=1.64$ m$^3$\,h$^{-1}$. The wetted radius grows with the nozzle size and the pressure; an empirical form is $R=3.4\sqrt{d\,P}$ with $d$ in mm and $P$ in bar, which gives $R=13.9$ m here.

\section{Spacing and application rate}
With the sprinklers on a rectangular spacing $S_l\times S_m$ the application rate is $I=q/(S_lS_m)$; here $I=7.6$ mm\,h$^{-1}$, which must be below the infiltration rate of the soil, or runoff follows. The spacing is chosen at 50 to 65 per cent of the wetted diameter, less in wind.

\section{Uniformity}
Christiansen's coefficient is $CU=100\bigl(1-\sum|z_i-\bar z|/\sum z_i\bigr)$ from catch-can depths $z_i$; 85 per cent or more is the target. For the depths 11, 12, 9, 13, 10, 12, 8, 11 mm: mean 10.75, $\sum|z-\bar z|=10.50$, $CU=87.8$ per cent.

